\begin{longtable}{
		|p{\dimexpr 0.3\textwidth - 11.28pt\relax}
		|p{\dimexpr 0.3667\textwidth - 13.79pt\relax}
		|p{\dimexpr 0.3333\textwidth - 12.53pt\relax}|
	}
	\caption{Элементы вкладки} \\
	\hline
	\textbf{Элемент} & \textbf{Значение} & \textbf{Назначение} \\
	\hline
	\endfirsthead
	\multicolumn{3}{l}{\textit{Продолжение таблицы~3.2}} \\
	\hline
	\textbf{Элемент} & \textbf{Значение} & \textbf{Назначение} \\
	\hline
	\endhead
	Поле ``Канал UART'' & 0--15 & Задает идентификатор канала, используемый при вызове \texttt{SendData}. Значение по умолчанию --- 0. \\
	\hline
	Поле ``Байты'' & Последовательность байтов & Используется для ручного ввода данных в шестнадцатеричном виде. \\
	\hline
	Поле ``Файл'' & Путь к выбранному файлу & Отображает путь к файлу, содержимое которого будет отправлено вместо ручного ввода. Поле недоступно для непосредственного редактирования. \\
	\hline
	Кнопка ``Обзор...'' & --- & Открывает окно выбора файла с данными. \\
	\hline
	Кнопка ``Отправить'' & --- & Выполняет отправку вручную введенных байтов либо содержимого выбранного файла. \\
	\hline
	Кнопка ``Запустить подписку'' & --- & Открывает потоковый вызов \texttt{Subscribe}. Во время активной подписки кнопка имеет текст ``Остановить подписку''. \\
	\hline
	Поле ``Отправленные данные'' & --- & Содержит сведения о выполненных операциях отправки. \\
	\hline
	Поле ``Принятые данные'' & --- & Содержит принятые пакеты, сообщения об ошибках и события управления подпиской. \\
	\hline
	Строка состояния & --- & Показывает состояние подключения, отправки, выбора файла, подписки и получения данных. \\
	\hline
\end{longtable}

Подключение выполняется автоматически во время создания вкладки. Графический
интерфейс загружает адрес микросервиса из конфигурационного файла и создает
gRPC-клиентское соединение.

До завершения попытки подключения поля канала, байтов и файла, а также кнопки
``Обзор...'', ``Отправить'' и ``Запустить подписку'' заблокированы.

При успешном подключении элементы управления разблокируются, а строка
состояния содержит сообщение:

\begin{center}
	\texttt{Статус: подключено к микросервису <адрес>}
\end{center}

Успешное соединение графического интерфейса с микросервисом не гарантирует
доступность драйвера RS-485. Ошибка соединения между микросервисом и драйвером
может быть получена при вызове \texttt{SendData} или во время работы
\texttt{Subscribe}.

Для ручной отправки данных необходимо:

\begin{enumerate}
	\item выбрать значение в поле ``Канал UART'';
	\item убедиться, что поле ``Файл'' пусто;
	\item ввести последовательность байтов в поле ``Байты'';
	\item нажать кнопку ``Отправить''.
\end{enumerate}

Байты задаются в шестнадцатеричном виде и разделяются пробелами. Каждый байт
должен находиться в диапазоне \texttt{00--FF}.

Пример корректного ввода:

\begin{center}
	\texttt{01 02 0A FF}
\end{center}

Регистр шестнадцатеричных цифр значения не имеет. Например, значения
\texttt{0a} и \texttt{0A} обозначают один и тот же байт.

При редактировании поля ``Байты'' ранее выбранный путь к файлу автоматически
очищается. После этого источником данных становится ручной ввод.

Для отправки содержимого файла необходимо:

\begin{enumerate}
	\item выбрать значение в поле ``Канал UART'';
	\item нажать кнопку ``Обзор...'';
	\item выбрать требуемый файл;
	\item убедиться, что путь появился в поле ``Файл'';
	\item нажать кнопку ``Отправить''.
\end{enumerate}

После выбора файла поле ``Байты'' автоматически очищается. Если поле ``Файл''
не пусто, при отправке используется содержимое файла, а не ручной ввод.

Файл передается как последовательность байтов. В журнале отправленных данных
источник отмечается как \texttt{file: <путь>}, а данные обозначаются строкой
\texttt{binary file}. При ручном вводе источник обозначается как
\texttt{manual input}.

Если оператор закрыл окно выбора файла, не выбрав файл, текущее состояние
полей не изменяется.

При выполнении команды $\mathtt{SendData}$ после нажатия кнопки ``Отправить'' интерфейс формирует запрос с идентификатором
канала и данными. Затем выполняется gRPC-вызов \texttt{SendData}.

В поле ``Отправленные данные'' записываются:

\begin{itemize}
	\item \texttt{channel\_id} --- идентификатор канала;
	\item \texttt{source} --- источник данных: ручной ввод или файл;
	\item \texttt{data} --- введенные байты либо обозначение двоичного файла;
	\item \texttt{success} --- признак успешного выполнения;
	\item \texttt{message} --- сообщение микросервиса;
	\item \texttt{exception} --- текст исключения, если операция не была выполнена.
\end{itemize}

При выполнении команды $\mathtt{Subscribe}$ для начала получения данных необходимо нажать кнопку
``Запустить подписку''. Интерфейс открывает потоковый gRPC-вызов
\texttt{Subscribe}, после чего:

\begin{itemize}
	\item кнопка меняет текст на ``Остановить подписку'';
	\item в поле ``Принятые данные'' появляется сообщение \texttt{Subscribe started. Waiting for data...};
	\item строка состояния сообщает о запуске подписки.
\end{itemize}

Для каждого результата, полученного из потока, в поле ``Принятые данные''
выводятся:

\begin{itemize}
	\item \texttt{channel\_id} --- идентификатор канала;
	\item \texttt{success} --- признак успешного получения;
	\item \texttt{data} --- принятые байты в шестнадцатеричном виде;
	\item \texttt{message} --- сообщение об ошибке либо пустая строка.
\end{itemize}

Для остановки активной подписки оператор нажимает кнопку
``Остановить подписку''. После этого:

\begin{itemize}
	\item потоковый вызов отменяется;
	\item кнопка возвращает текст ``Запустить подписку'';
	\item в поле ``Принятые данные'' добавляется сообщение \texttt{Subscribe stopped.};
	\item строка состояния сообщает об остановке подписки.
\end{itemize}